\uuid{pbyD}
\titre{Théorème de Bézout}
\auteur{bodin ; croizat ; sacre}
\organisation{exo7}
\licence{CC BY-NC-SA 4.0}
\datecreate{2022-03-08}
\niveau{L1}
\module{Algèbre}
\chapitre{Arithmétique dans Z}
\sousChapitre{Nombres premiers, nombres premiers entre eux}
\difficulte{2}
\theme{Théorème de Bézout}
\exo7qcmid{hilisit.arithmetique.018}

\contenu{
\begin{question}
\qclassification{103.04}
    Soient deux entiers positifs $a,b$, on calcule le pgcd de $a$ et $b$ par l'algorithme d'Euclide.
    La première étape est d'écrire la division euclidienne de $a$ par $b$ : $a=bq+r$.
    Quelle est la second étape ?
    \begin{answers}
        \bad{La division de $a$ par $r$.}
        \good{La division de $b$ par $r$.}
        \bad{La division de $q$ par $r$.}
        \bad{Cela dépend des valeurs de $a$ et $b$.}
    \end{answers}
    \begin{explanations}
        Une conséquence de l’égalité est que $\pgcd(a,b)=\pgcd(b,r)$. On remplace donc $a$ par $b$ et $b$ par $r$ dans l’étape suivante, c’est-à-dire qu’on fait la division euclidienne de $b$ par $r$.
    \end{explanations}
\end{question}
}
